\section{Specification and Verification Boundary}
\label{sec:specification}

The overview deliberately hid most notation.  This section states the
mathematical model, the representative operation condition, and the
caller-visible form of the running contract.

We use IEEE~754 binary64 with round-to-nearest, ties-to-even.  Let
$u=2^{-53}$ be unit roundoff, $\eta=2^{-1074}$ the smallest positive
subnormal value, and $\Omega=2^{1024}$ the power of two immediately above the
largest finite binary64 binade.  Flocq supplies the floating-point semantics
and rounding definitions used in the development~\cite{Boldo2011}.

\begin{definition}[Representative operation condition]
\label[definition]{def:dw-times-fp-cond}
The predicate $\DWTimesFPCond(x_h,x_\ell,y)$ combines the hypotheses used by the
bounded-format theorem for the running example.  At the paper level, its main
clauses state that the relevant values are finite, $(x_h,x_\ell)$ satisfies
\eqref{eq:dw-representation}, and the leading product has sufficient upper-range
headroom for the intermediate computation.
\end{definition}

The development's condition is stated over Flocq values and their real
interpretations.  The C check cannot in general test an exact real expression
directly.  Instead, its VST proof shows that the Boolean tests made by the
program imply the condition.  This keeps the numerical condition independent
of a particular implementation of its sufficient checker.

The principal result for the running example can be stated without exposing
the proof-assistant representation of every connective.

\begin{theorem}[Guarded DW$\times$FP contract]
\label{thm:guarded-dw-times-fp}
Let $x$ point to a stored binary64 pair $(x_h,x_\ell)$, and let $z$ point to
valid, suitably separated result storage.  If the guarded DW$\times$FP call
returns normally with $(z_h,z_\ell)$, the returned pair is the result of the
binary64 functional model, satisfies the error bound stated in
\eqref{eq:dw-times-fp-mixed}, and satisfies the operation's output representation
property.  If the check fails, the wrapper follows its specified nonreturning
error path without invoking the numerical kernel.
\end{theorem}

The formal VST contract divides these facts among three zones.  Pure numerical
propositions belong in \texttt{PROP}; function arguments and return values are
described in \texttt{PARAMS}/\texttt{RETURN}; and ownership of the input and
output structures belongs in \texttt{SEP}.  Write $R(x)$ for the memory
representation, $C(x)$ for the numerical condition, and $G(x,z)$ for the
operation result.  Schematically, the three relevant contracts are
\[
\begin{array}{c}
\{R(x)\}\ \texttt{check}(x)\
  \ \{b=\mathsf{true}\Rightarrow C(x)\},\\[2pt]
\{R(x)\land C(x)\}\ \texttt{kernel}(x)\ \{R(z)\land G(x,z)\},\\[2pt]
\{R(x)\}\ \texttt{guarded}(x)\
  \ \{\text{normal return}\Rightarrow G(x,z)\}.
\end{array}
\]
The logical conjunctions in this schematic display should not be confused with
separating conjunction between memory resources.

The theorem concerns functional correctness and numerical accuracy of normally
returning calls.  It is not a termination theorem for arbitrary environments,
a complexity result, or a claim that the pair is a correctly rounded binary128
result.  Its compilation-level interpretation is limited to the verified path
and assumptions described in Section~\ref{sec:overview-pipeline}.
